\documentclass{article}

\textwidth=6.5in
\textheight=8.5in
\voffset=-54pt
\oddsidemargin=0in
\evensidemargin=0in
\setlength{\parskip}{8pt}
\setlength{\parindent}{0pt}

\usepackage{amsmath, amssymb}
\usepackage{ifthen}

\usepackage[inline]{enumitem}
\setlist{topsep=0pt, parsep=8pt}

\usepackage[rgb,table]{xcolor}\convertcolorsUfalse
\usepackage{graphicx}

% change hyperref options
\PassOptionsToPackage{hyphens}{url}
\usepackage[bookmarks,colorlinks,linkcolor=black,citecolor=blue,pdfstartview=FitH,urlcolor=blue]{hyperref}
\hypersetup{pdfpagemode=UseNone}
\usepackage{bookmark}

\usepackage{graphicx}
\usepackage{multicol}
\usepackage{framed}
\usepackage{array}

\usepackage{icomma}

\usepackage{soul}

\usepackage{pgfplots}
\pgfplotsset{compat=1.18}  
\pgfplotscreateplotcyclelist{mycolor}{black, blue, red, green!70!black, magenta, purple, cyan, orange }  
\pgfplotsset{
  disabledatascaling,
  cycle list=tikzcolor, 
  axis lines=middle,
  every axis plot/.style={no marks, thick},
  every axis/.style={
    enlarge x limits=false,
    enlarge y limits=auto,
    xlabel style={at={(current axis.right of origin)},anchor=west},
    ylabel style={at={(current axis.above origin)},anchor=south},
    % extra x and y ticks for asymptotes
    extra x tick style={grid=major, grid style={dashed,black}},
    extra y tick style={grid=major, grid style={dashed,black}},
    /pgf/number format/1000 sep={},
  },    
  tick label style = {font=\footnotesize, fill=\currentbgcolor, inner sep=0pt},    
  xticklabel shift=1pt,    
  scaled ticks=false,
  scaled y ticks=false,
  unbounded coords=jump,
  samples=101,
  minor tick length=0,
  scatter/classes={
    open={mark=*,draw=black,fill=white, mark size=2, thin},
    closed={mark=*,draw=black,fill=black, mark size=2, thin}
  },
  width=3in,
}
\pgfplotsset{open/.style={only marks, mark=*, fill=white, thin, mark size=2}}
\pgfplotsset{closed/.style={only marks, mark=*, draw=black, fill=black,
    thin, mark size=2}}
\pgfplotsset{labeled/.style={nodes near coords, only marks, mark=*, draw=black, fill=black, thin, mark size=2, point meta=explicit symbolic}}

\usepackage{tikz}
\usetikzlibrary{
  matrix,%
  % enables the snake paths
  decorations.pathmorphing,%
  % enables bigger arrows
  decorations.markings,%
  decorations.pathreplacing,%
  decorations.text,%
  calc,%
  patterns,%
  shapes.geometric,%
  arrows,
  decorations.shapes,
  positioning,
  plotmarks,
  intersections
}
\tikzset{>=latex} % sets default arrow type in tikz
\tikzset{font=\normalsize}
\tikzset{mark size=1.5pt, mark options=thin}
\tikzset{pin distance=4pt,
  every pin edge/.style={<-, thin, shorten <= -2pt}}

\interfootnotelinepenalty=10000
\renewcommand{\thefootnote}{(\arabic{footnote})}

\usepackage{multirow, bigdelim}

% change to \usepackage[sols]{handout} to see solutions
\usepackage[sols, notes, workshop]{handout}
\renewcommand{\workshopnote}{CA Notes.}    

\newcommand\iscurrentenumi[1]{%
  \ifthenelse{\equal{#1}{\theenumi}}{}{#1}}
\setenumerate[1]{ref=\#\arabic*}
\setenumerate[2]{ref=\protect\iscurrentenumi{\theenumi}(\alph*)}
\newcommand\enumiiprefix[2]{%
  \ifthenelse{{\equal{#1}{\theenumi}}\AND{\equal{#2}{(\alph{enumii})}}}{}{}%
  \ifthenelse{{\equal{#1}{\theenumi}}\AND{\NOT{\equal{#2}{(\alph{enumii})}}}}{#2}{}%
  \ifthenelse{\equal{#1}{\theenumi}}{}{#1#2}%
}
\setenumerate[3]{ref=\protect\enumiiprefix{\theenumi}{(\alph{enumii})}\roman*}

\definecolor{purple}{rgb}{0.5,0,1.0}
\definecolor{hotpink}{rgb}{1, 0, 0.5}
\definecolor{darkorange}{rgb}{1, 0.4, 0}

\definecolor{skyblue}{rgb}{0, 0.5, 1}

\usepackage{fancyhdr}
\lhead{\textcolor{white}{This content is protected and may not be shared, uploaded, or distributed.}}
\rhead{}
\rfoot{\vspace*{\fill}\footnotesize\textcopyright{} \emph{Harvard Math Ma}}
\renewcommand{\headrulewidth}{0pt}
\pagestyle{fancy}
\begin{document}

\htitle{Workshop 9}

\begin{workshop}
  Today's problems focus on practice with exponentials and logs, which will hopefully help students prepare for the upcoming skills check.

  \#4 previews solving equations with logarithms and exponents (Friday's lesson). Please do \# 1-3, 5 first and then have students go back and attempt \#4 if there is extra time.
\end{workshop}

\begin{enumerate}
\item  A population of mosquitoes is growing exponentially.  When there are 5000 mosquitoes, the mosquito population is growing at a rate of 100 mosquitoes per month. Let $f(t)$ be the mosquito population at time $t$, where $t$ is measured in months. 

  \begin{enumerate}

  \item
    \begin{problem}[\vfill]
      Since the mosquito population is growing exponentially, the population is changing at a rate proportional to its size. What is the constant of proportionality in this case? In other words, fill in the blank: \probonly{$\vphantom{\boxed{\int_1^2}}$}$f'(t) = \fitb{0.25in}{}\ f(t)$.
    \end{problem}

    \begin{solution} 
      $$f'(t) = kf(t)$$ 
      When $f(t)=5000$, we have that $f'(t) = 100$, so then
      $$100=k \cdot 5000 $$ 
      $$k = \frac{1}{50}$$  
    \end{solution} 

  \item
    \begin{problem}[\vfill]
      How quickly is the population growing when there are 6000 mosquitoes? Be sure to include units.
    \end{problem}

    \begin{solution}
      $f'(t)=0.02(6000)=\frac{1}{50}6000=120$ mosquitoes per month.
    \end{solution}

  \item % added 2023

    \begin{workshop}
      An important point here is that, when we interpret $f^{-1}(t)$, $t$ no longer represents time! Instead, it represents a number of mosquitoes. 
    \end{workshop}

    \begin{problem}[\vfill]
      Is $f(t)$ invertible? How do you know? If so, interpret the meaning of $f^{-1}(t)$. 
    \end{problem}

    \begin{solution}
        It is invertible, because $f(t)$ is strictly increasing. The inverse function, $f^{-1}$ outputs the time needed for the population to reach a certain size.  
    \end{solution}

  \end{enumerate}
  Suppose in addition that, when $t = 0$, there are 3000 mosquitoes. 
  \begin{enumerate}[resume]
  \item

    \begin{problem}[\stretch{2}]
      Find a formula for $f(t)$.

      \emph{Hint:} Since we're told that $f(t)$ is exponential, we can write it as $f(t) = C e^{kt}$ or as $f(t) = C \cdot b^t$. What's $C$? Can you use your answer to (a) to find either $b$ or $k$? 
    \end{problem}

    \begin{solution}
        We know that:
        \begin{itemize}
            \item $f(t) = Ce^{kt}$
            \item $f(0) = Ce^{k(0)} = C = 3000$
            \item $f'(t) = k \cdot 3000 \cdot e^{kt}$ (because $\dfrac{d}{dx}\Big[e^{kx}\Big]=ke^{kx}$)
        \end{itemize}
        And since $f'(t) = k \cdot 3000 \cdot e^{kt}$, and $f'(t) - \dfrac{1}{50} \cdot f(t)$ (from Part (a)), $k$ must equal $\dfrac{1}{50}$ and thus 
        $$f(t) = 3000 e^{\frac{1}{50}t}$$
    \end{solution}

  \item % added 2023

    \begin{problem}[\nstretch{1.5}]
      How long does it take the mosquito population to double? Give an exact answer.
    \end{problem}

    \begin{solution}
        To answer this, we need to solve:
        \begin{align*}
            3000 e^{\frac{1}{50}t} = 6000 \\
            e^{\frac{1}{50}t} = 2 \\
            \ln(e^{\frac{1}{50}t}) = \ln(2)\\
            \dfrac{1}{50} t = \ln(2) \\
            t = 50 \ln(2)\\
        \end{align*}
        
        
    \end{solution}

  \item % added 2023

    \begin{problem}[\stretch{3}]
      You should have found that $f$ is invertible. Find a formula for $f^{-1}(t)$. 
    \end{problem}

    \begin{solution}
        \begin{align*}
        \intertext{Set $f(t) = y$} \\
            y = 3000 e^{\frac{1}{50}t} \\
            \intertext{Swap $t$ and $y$} \\
           t = 3000 e^{\frac{1}{50}y} \\
            \intertext{Solve for $y$} \\
            \ln(t) = \ln(3000) + \dfrac{1}{50} y \cdot \ln(e) \\
            \ln(t) - \ln(3000) = \dfrac{1}{50} y \\
            50 \ln\Big(\frac{t}{3000}\Big) = y = f^{-1}(t)\\
        \end{align*}
    \end{solution}

  \end{enumerate}

\item 

  \begin{workshop}
    Practice with log rules and graph transformations
  \end{workshop}

  In each part, you're given a function defined in terms of logs. Simplify the formula for the function as much as possible, and then pick its graph (choices 1 - 5). Explain how you made your choice. 

  \begin{center}
    \begin{tikzpicture}
      \begin{axis}[xtick={1}, ymin=-6, ytick={1}, clip=false]
        \addplot+[domain=0:6, green!80!black, dashdotted] {ln(6)} node[right]{(3)}; % 
        \addplot+[domain=0:6, skyblue, densely dotted] {ln(2)} node[right]{(4)}; %
        \addplot+[domain=0:6, hotpink, restrict y to domain=-6:10] {7/3*ln(x) / ln(2)} node[right]{(1)}; %
        \addplot+[domain=1/729:6, samples=201, purple] {-ln(x)/ln(3)} node[right]{(5)}; %
        \addplot+[domain=0:6, samples=201, darkorange, densely dashed] {ln(5*x^2)/ln(3)} node[right]{(2)}; %
      \end{axis}
    \end{tikzpicture}
  \end{center}

  \begin{enumerate}
  \item

    \begin{workshop}
      This simplifies to just $\ln 2$, which is a constant. So, the graph is either (3) or (4). To decide which, students need to know whether $\ln 2$ is greater than or less than 1.
    \end{workshop}

    \begin{problem}[\vfill]
      $f(x) = \ln(2x) - \ln(x)$
    \end{problem}

    \begin{solution}
      % ln (2) is a constant < 1
           $\ln(2x)-\ln(x) = \ln\Big(\dfrac{2x}{x}\Big) = \ln(2)$. 

    This is a constant function, so it is graph (4).
    \end{solution}

  \item
    \begin{problem}[\vfill]
      $g(x) = \log_2 \left( \sqrt[3]{x} \right) + \log_2 \left( 8 x^3 \right) - \log_2 (8x)$
    \end{problem}

    \begin{solution}
      % vertical stretch with no flip, so still has x-intercept 1
    \begin{align*}
    &  g(x) = \log_2 \left( \sqrt[3]{x} \right) + \log_2 \left( 8 x^3 \right) - \log_2 (8x)  \\
      & =\dfrac{1}{3}\log_2(x) + \log_2(8) + 3\log_2(x) - \log_2(8) - \log_2(x)\\
        & = \frac{5}{3}\log_2(x) \\
    \end{align*}
    This is graph (1). 
    \end{solution}

  \item
    \begin{problem}[\vfill]
      $h(x) = \log_3 (7) - \log_3 (7x)$
    \end{problem}

    \begin{solution}
      % vertical flip of \log_3 x

      $h(x) = \log_3 (7) - \log_3 (7x) = \log_3\Big(\dfrac{7}{7x}\Big) = \log_3\Big(\dfrac{1}{x}\Big) = \log_3(x^{-1}) = -\log_3(x)$
      
      This is graph (5).
    \end{solution}

  \item
    \begin{problem}[\vfill\newpage]
      $j(x) = \log_3 (5x^2)$
    \end{problem}

    \begin{solution}
      % vertical stretch and shift up of \log_3 x, so x-intercept no longer at 1 
        $j(x) = \log_3 (5x^2) = \log_3(5) + \log_3(x^2) = \log_3(5)+2\log_3(x)$.

        This is graph (2).
    \end{solution}

  \end{enumerate}

\item  

  Let $\displaystyle f(x)=\frac{\ln \left(\sqrt{3x} \right)}{2}+3$.

  \begin{enumerate}
  \item

    \begin{problem}[0.5in]
      What's the domain of $f$?
    \end{problem}

    \begin{solution}
        The domain of $f$ is $(0, \infty)$. 
    \end{solution}

  \item
    \begin{problem}[1in]
      Find $f'(x)$.
    \end{problem}

        \begin{solution}
        It's helpful to note that $f$ can be rewritten as:
        \begin{align*}
            f(x)=\frac{\ln \left(\sqrt{3x} \right)}{2}+3 \\
            = \frac{\ln(3x)^{\frac{1}{2}}}{2} + 3 \\
             = \frac{\frac{1}{2}\ln(3x)}{2} + 3 \\
             = \frac{\ln(3x)}{4 } +3 \\
            = \frac{\ln(3)+\ln(x)}{4} + 3\\
            = \frac{\ln(3)}{4} + \frac{\ln(x)}{4} + 3\\
        \end{align*}

        Then, we can see that $f'(x) = \dfrac{1}{4x}$
    \end{solution}


  \item \label{gottlieb-14.1.8-increasing}

    \begin{grading}
      2 points: 1 for calculating $f'$ correctly, 1 for using that to show  that $f$ is increasing everywhere
    \end{grading}



    \begin{problem}[\vfill]
      Where is $f$ increasing? Where is $f$ decreasing? 
    \end{problem}

    \begin{solution}
      \begin{align*}
        f(x) &= \frac{\frac{1}{2} \ln (3x)}{2} + 3 \\
        &= \frac{\frac{1}{2} (\ln 3 + \ln x)}{2} + 3 \\
        &= \frac{1}{4} ln x + \frac{\ln 3}{4} + 3, \\
        \shortintertext{so}
        f'(x) &= \frac{1}{4} \cdot \frac{1}{x} \\
        &= \frac{1}{4x}
      \end{align*}
      This is positive everywhere on $(0, \infty)$, so $f$ is increasing on $(0, \infty)$. 
    \end{solution} 

  \item

    \begin{grading}
      1 point
    \end{grading}

    \begin{problem}[0.75in]
      Explain why your answer to \ref{gottlieb-14.1.8-increasing} shows that $f$ is invertible. 
    \end{problem}

    \begin{solution} 
      Since the function is always increasing, it will pass the horizontal test and be invertible.  
    \end{solution}

  \item

    \begin{grading}
      2 points: 0.5 for knowing to switch $x$ and $y$, 1.5 for correctly solving for $y$ (give partial credit if they make partial progress)
    \end{grading}

    \begin{problem}[\vfill\newpage]
      Find a formula for $f^{-1}(x)$.
    \end{problem}

    \begin{solution}
      We start with $y = \frac{\ln(\sqrt{3y})}{2}+3$, switch $x$ and $y$ (since an inverse function switches the input and output), and solve for $y$:
      \begin{align*}
        x &= \frac{\ln(\sqrt{3y})}{2}+3 \\
        x-3  &=  \frac{\ln(\sqrt{3y})}{2} \\
        2(x-3)  &=  \ln(\sqrt{3y})  \\
        e^{2(x-3)} &= \sqrt{3y} \\
        e^{4(x-3)}  &=  3y \\
        \boxed{\frac{e^{4(x-3) } }{3}} &= y
      \end{align*} 
    \end{solution}

  \end{enumerate}

\item  Solve for $x$.
  \begin{enumerate}
  \item
    \begin{problem}[\vfill]
      $ \log_{10} x + \log_{10} (2x) + \log_{10} 5 = 55$ (Please simplify your answer.) 
    \end{problem}

    \begin{solution}
      $ x=10^{27}$ (take 10 to the power of each side, and careful with the exponent rules!)
    \end{solution}

  \item
    \begin{problem}[\vfill]
      $ 3^{2x} - 10 \cdot 3^{x} + 24 = 0 $ (\emph{Hint:} Substitute $u = 3^x$. How can you express $3^{2x}$ in terms of $u$?) 
    \end{problem}

    \begin{solution}
      $x=\log_3 4, \; \log_3 6 $ (let $y=3^x$ and solve the quadratic in $y$ first).
    \end{solution}

  \item
    \begin{problem}[\vfill\newpage]
      $2 \log_2 (x+1) - \log_2 (x+2) = -1 $
    \end{problem}

    \begin{solution}
      $x=0$
    \end{solution}

  \end{enumerate}

\item 
  Find the derivative of each of the following functions. 

  \begin{enumerate}

  \item
    \begin{problem}[\stretch{0.25}]
      $f(x)=\ln \left(\dfrac{5}{x} \right)$
    \end{problem}

    \begin{solution}
      $f(x)=\ln5-\ln x$, so $f'(x)=-\frac{1}{x}$
    \end{solution}

  \item
    \begin{problem}[\stretch{0.25}]
      $f(x)= \dfrac{\ln \left(7x^3e^{5x} \right)}{e^2}$
    \end{problem}

    \begin{solution}
        \begin{align*}
            f(x) &= \dfrac{\ln \left(7x^3e^{5x} \right)}{e^2} \\
             &= \frac{\ln(7) + \ln(x^3) + \ln(e^{5x})}{e^2} \\
             &= \frac{1}{e^2}\Big(\ln(7)+ 3 \ln(x) + 5x\Big) \\
    \intertext{So then we have}
    f'(x) &= \frac{1}{e^2}\Big(\frac{3}{x} + 5\Big)
        \end{align*}
    \end{solution}

  \item
    \begin{problem}[\stretch{0.25}]
      $j(w) = 2^w \cdot 3^{-w/2}$
    \end{problem}

    \begin{solution}
        \begin{align*}
            j(w) &= 2^w \cdot 3^{-w/2} \\
            & = 2^w \cdot (3^{-1/2})^w \\
            & = \Big(\dfrac{2}{3^{1/2}}\Big)^w \\
        \intertext{So then we have:} 
        & j'(w) = \Big(\dfrac{2}{3^{1/2}}\Big)^w \cdot \ln\Big(\frac{2}{3^{1/2}}\Big) \\
        \end{align*}
    \end{solution}

  \item
    \begin{problem}[\stretch{0.25}]
      $h(t) = e^{2t + 3 \ln t}$
    \end{problem}

    \begin{solution}
        \begin{align*}
            h(t) &= e^{2t + 3 \ln t} \\
            & = e^{2t}e^{3\ln(t)} \\
            & = e^{2t}e^{\ln(t)^3} \\
            & = e^{2t}t^3 \\
            \intertext{So then we have:} 
           & h'(t) = 2e^{2t}t^3 + 3e^{2t}t^2 \\
        \end{align*}
    \end{solution}

  \end{enumerate}

\end{enumerate}
\end{document}
